\chapter{The Container}

Three rules govern reading an APFS container: block zero locates,
it does not authorise; checkpoint selection picks the highest
checksum-valid candidate; and the selected state must not advance
during the scan. This chapter develops each rule.

\section{Block Zero As Locator}

APFS is copy-on-write at the container level. Each transaction writes a
new container superblock (\texttt{nx\_superblock\_t}, henceforth NXSB)
into the checkpoint descriptor area. Block zero is the stable known
location at which a reader can find the descriptor area. It is updated
lazily; the authoritative live NXSB is one of the copies in the
descriptor ring.

The pieces of block zero a reader needs:

\begin{itemize}
  \item \texttt{block\_size} at offset \texttt{0x24}: validated to be a
    power of two within \texttt{[4096, 65536]}.
  \item \texttt{nx\_xp\_desc\_blocks} at offset \texttt{0x68}: the
    length of the descriptor area in blocks.
  \item \texttt{nx\_xp\_desc\_base} at offset \texttt{0x70}: the first
    block of the descriptor area. The high bit signals a
    \emph{non-contiguous} layout, which this work does not support
    and which triggers a typed hard stop. A non-contiguous descriptor area
    would require a separate checkpoint-mapping-tree implementation.
\end{itemize}

Block zero is itself an APFS object and must validate as one: NXSB
magic at offset \texttt{0x20}, object type \texttt{NX\_SUPERBLOCK},
object id equal to the constant NX-superblock OID (1), and a valid
Fletcher-64 checksum. A block zero that does not satisfy these is not
an APFS container.

\textbf{Observation.} The \texttt{block\_size} and descriptor base in
block zero must agree with the selected NXSB's view of the same
fields. A discrepancy means block zero is stale relative to a
descriptor write that did not (yet) propagate. The parser rejects
the source rather than guessing which view is correct.

\begin{figure}[h]
\centering
\begin{tikzpicture}[
  every node/.style={font=\small},
  block/.style={draw, minimum width=1.05cm, minimum height=0.8cm,
                align=center, font=\footnotesize},
  locator/.style={block, fill=gray!12},
  nxsb/.style={block, fill=blue!10},
  nxsbsel/.style={block, fill=blue!30, very thick},
  map/.style={block, fill=orange!18},
  emptyc/.style={block, fill=white, dashed},
  arrow/.style={-Stealth, thick},
]
  \node[locator, minimum width=1.6cm, minimum height=1.5cm] (b0)
    {block 0\\\scriptsize locator};

  \node[nxsb, right=1.8cm of b0, yshift=0.2cm] (n1) {NXSB\\\scriptsize xid 18};
  \node[nxsb, right=0.05cm of n1] (n2) {NXSB\\\scriptsize xid 19};
  \node[map,  right=0.05cm of n2] (m1) {map};
  \node[nxsb, right=0.05cm of m1] (n3) {NXSB\\\scriptsize xid 20};
  \node[nxsbsel, right=0.05cm of n3] (n4) {NXSB\\\scriptsize xid \textbf{21}};
  \node[map,  right=0.05cm of n4] (m2) {map\\\scriptsize last};
  \node[emptyc, right=0.05cm of m2] (e1) {};
  \node[emptyc, right=0.05cm of e1] (e2) {};

  \node[fit=(n1)(e2), draw=none, label={[font=\footnotesize, yshift=-0.05cm]
        above:checkpoint descriptor area}] (ring) {};

  \draw[arrow] (b0.east) -- ++(0.4,0) |- (ring.west)
    node[midway, above, font=\scriptsize] {desc base};

  \draw[arrow, dashed]
    (n4.south) ++(0,-0.1) -- ++(0,-0.7)
    node[below, font=\scriptsize, align=center]
      {selected NXSB:\\highest checksum-valid xid};
\end{tikzpicture}
\caption{Block zero locates the descriptor ring; the selected NXSB is
  the highest checksum-valid candidate in the ring. Several earlier
  NXSBs may remain valid; \texttt{map} blocks carry checkpoint-map
  contents; trailing slots may be unwritten between commits.}
\end{figure}

\section{Selecting The Checkpoint}

The descriptor area is a contiguous run of \texttt{nx\_xp\_desc\_blocks}
blocks. Some are checkpoint maps (chapter~\ref{ssec:checkpoint-map});
some are NXSB candidates. The selection rule:

\begin{enumerate}
  \item Iterate every block in the descriptor range.
  \item For each, check the NXSB magic at offset \texttt{0x20}.
  \item If the magic matches, validate object type
    (\texttt{NX\_SUPERBLOCK}) and Fletcher-64 checksum.
  \item Record the candidate.
  \item After the scan, pick the candidate with the highest XID.
\end{enumerate}

The highest valid XID is the selected NXSB. Its XID becomes the
\emph{scan XID}, and every subsequent object resolution for the run is
performed at this XID.

Blocks that carry NXSB magic but fail object-type or checksum
validation are recorded as skipped descriptors with the reason. They
are not silently ignored; a future support broadening must justify
why they were skipped.

\textbf{Observation.} Three or four NXSB candidates is typical for a
recently written container. The first execution of EX-15 against the
EX-14 fixture saw four candidates with XIDs 18--21; all four passed
strict validation. Multiple candidates is not a sign of corruption.
Zero candidates is.

\subsection{Why The Highest Valid XID}

The candidates are not chronologically ordered in the descriptor area;
they are written into a ring. The XID is the only monotonic identity
for transaction order. The highest checksum-valid XID is the most
recent transaction the writer completed. A higher XID with an invalid
checksum is a transaction the writer started but did not finish; the
parser treats those as failed and ignores them, with the result that
selection is always self-consistent.

The rule has a corollary: a parser that picks a lower XID when the
highest is valid has chosen a stale state for no reason and will
disagree with anything else that picks the latest.

\subsection{Live State Is Not Pinned}

\textbf{Observation.} On a mounted APFS volume that other processes
are using, the highest valid XID advances during the scan. Two
reads of ``latest'' may return different transactions; mixing
objects across transactions produces an incoherent namespace. EX-01
demonstrated this on a mounted lab image under write churn.

Raw mode therefore only runs on sources where state can be pinned:
a detached image (the \texttt{.dmg} layer fixes the visible state
at detach time), or a caller-supplied raw device under an explicit
commitment that the checkpoint will not advance. Live mounted
volumes are readable but not supported; chapter~11 develops the
vocabulary.

\section{The Checkpoint Map Ring}
\label{ssec:checkpoint-map}

A selected NXSB references a chain of checkpoint maps
(\texttt{checkpoint\_map\_phys\_t}) in the descriptor area. The map
chain carries the ephemeral object identifier table for the selected
checkpoint: the OIDs and physical addresses of objects that live in
memory between transactions and are written into the checkpoint data
area at commit. The reader needs the map to resolve those ephemerals.

The walk rule:

\begin{enumerate}
  \item Start at \texttt{nx\_xp\_desc\_index} in the selected NXSB.
  \item Read the block; validate as an APFS object with type
    \texttt{OBJECT\_TYPE\_CHECKPOINT\_MAP}.
  \item Record the (\texttt{oid}, \texttt{paddr}) mappings.
  \item Advance through subsequent map blocks in the descriptor area.
  \item Stop when a block carries the \texttt{CHECKPOINT\_MAP\_LAST}
    flag.
\end{enumerate}

A chain that runs out of length before the last flag is malformed.
The parser fails closed.

\textbf{Empirical correction (EX-11).}
\texttt{OBJECT\_TYPE\_CHECKPOINT\_MAP} blocks carry the
\texttt{OBJ\_PHYSICAL} storage flag, not \texttt{OBJ\_EPHEMERAL},
despite holding the ephemeral-object lookup table. The Apple
reference reads naturally as ``the checkpoint-map block itself is
ephemeral'' --- it is not. The validator expects
\texttt{Physical} for these blocks. A synthetic malformed-map
regression test covers the constant.

\section{Roots Reachable From The Selected Checkpoint}

The selected NXSB names five things the reader needs:

\begin{enumerate}
  \item \texttt{nx\_omap\_oid}: the OID of the container OMAP.
  \item \texttt{nx\_fs\_oid[]}: an array of volume superblock OIDs.
  \item Feature, incompatible-feature, and read-only-compatible
    feature masks.
  \item Reaper and space-manager OIDs (not used by an indexer;
    their presence
    is required, their content is not).
  \item Container UUID (used for source identity in caches and
    diagnostics).
\end{enumerate}

The incompatible-feature mask is the gate for format drift
(chapter~11). Any bit the reader does not recognise stops the run.

Block zero, the descriptor area, and the checkpoint map together
establish a single coherent transaction context. The OMAP
machinery of the next chapter operates inside that context.
